\section{Results and Discussion}
\label{sec:results}

\subsection{Comparison with semantic segmentation}

Table~\ref{tab:main} places the default YOLO configuration alongside nine
baselines. Averaged over three seeds YOLO26n-seg reaches $0.9113$ Dice and
$3.809$\,mm HD95 at $2.7$\,M parameters, against $0.9121$/$4.081$\,mm for
TransUNet ($102$\,M). The axes differ: on HD95 it is best in the benchmark by
$0.271$\,mm, twice the seed floor; on Dice it is indistinguishable from the four
leading baselines, tested over the $50$ patients with seeds averaged. Seed~0
scored over frames would instead give $0.9151$ and eight of nine.

%============================================================================
% Y1 -- YOLO vs semantic-segmentation baselines (auto-generated)
%============================================================================
\begin{table*}[t]
\centering
\caption{CAMUS official test split (50 patients). Dice on the $256\times256$ grid, HD95 at native resolution in mm. The detector row is the mean over three seeds, not the best of them. Best per column in \textbf{bold}. Markers: Holm-corrected paired Wilcoxon over patients against that mean, $^{*}$ $p<0.05$, $^{\dagger}$ not significant.}
\label{tab:main}
\small
\setlength{\tabcolsep}{4pt}
\begin{tabular}{l l c c c}
\toprule
Model & Paradigm & Dice & HD95 (mm) & Params (M) \\
\midrule
TransUNet$^{\dagger}$      & Hybrid      & \textbf{0.9121} & 4.08 & 102.1 \\
YOLO26-seg (ours, mean of 3 seeds) & Detector    & 0.9113 & \textbf{3.81} & 2.7 \\
nnU-Net$^{\dagger}$        & CNN         & 0.9099 & 4.27 & 15.6 \\
UNet-V1$^{\dagger}$        & CNN         & 0.9085 & 4.62 & 31.0 \\
UNet-ResNet$^{\dagger}$    & CNN         & 0.9081 & 4.26 & 29.1 \\
UNet-V2$^{*}$              & CNN         & 0.8963 & 5.29 & 33.0 \\
Swin-UNet$^{*}$            & Transformer & 0.8882 & 5.08 & 41.8 \\
DeepLabV3+$^{*}$           & CNN         & 0.8602 & 5.71 & 40.3 \\
FPN-UNet$^{*}$             & CNN         & 0.8596 & 6.21 & 32.2 \\
DenseContextU-Net$^{*}$    & CNN         & 0.8505 & 8.29 & 2.2 \\
\bottomrule
\end{tabular}
\end{table*}


Two external comparisons follow on the same split. Against CGQR-Net, a
contour-guided boundary-aware method reporting per-structure Dice of
$93.95$/$88.09$/$89.15$ (mean $90.40$), our seed mean is
$93.86$/$87.91$/$91.62$ (mean $91.13$): ahead on the atrium and the
three-class mean, marginally behind on both ventricular classes. Against a multi-task framework reporting
$0.951$ for the LV cavity at ED, our LV-endocardial Dice at ED is $0.9523$. More
important than either is the dataset authors' own reference point: average
inter-observer Dice on CAMUS is $0.899$~\cite{leclerc2019camus}. Every model in
Table~\ref{tab:main} above $0.90$ already operates beyond the level at which two
cardiologists agree, which is how the small differences there should be read.

\subsection{Latency under a single protocol}

Published latencies for these baselines were measured on other hardware and
cannot be mixed with ours, so we re-time every model on one GPU with one
protocol (batch 1, CUDA-synchronised, median of 50 after 10 warmup).

The result is a property of the architecture family rather than a tuning
outcome: \textbf{YOLO is near-invariant to input resolution}, $15.7$\,ms at
$256$\,px and $15.5$\,ms at $640$\,px, while every encoder--decoder scales with
pixel count ($4$--$8\times$: nnU-Net $10.7\rightarrow78.5$\,ms, UNet-V1
$17.7\rightarrow119.3$\,ms). At matched $640$\,px YOLO is $1.5$--$23\times$
faster than every baseline; end-to-end inference is $11.8$\,ms. Neither FP16
($16.7$\,ms) nor ONNX export ($23.5$\,ms) improved on the PyTorch path.

\subsection{Ejection fraction}

EF requires pairing each patient's ED and ES predictions in both views. With the
official CAMUS biplane Simpson implementation at native resolution, YOLO reaches
$9.24\%$ MAE at $r{=}0.756$: within the baselines' range ($6.66$--$9.74\%$,
$r{=}0.640$--$0.827$) but not leading them, so we claim no EF advantage. The
ceiling matters more than the ranking -- the \emph{ground-truth} masks through
the same pipeline give $5.44\%$ MAE at $r{=}0.863$, the irreducible gap between
mask-derived and clinically recorded EF, which consumes most of every model's
error.

One implementation note with wider consequences: the official Simpson code takes
the first contour by position rather than the largest. Ground truth is
single-component so this never manifests, but predictions carry occasional
spurious blobs -- on one baseline $6\%$ of frames -- which sufficed to drive EF
correlation from $r{=}0.911$ to $r{=}0.030$. A largest-component filter is
mandatory when computing EF from predictions.

\subsection{Discussion}

\textbf{Nothing we varied is the bottleneck.} A perfect model under our
decomposition scores $0.997$ Dice at one-pixel HD95, yet no configuration change
improves on the stock setting beyond the $0.134$\,mm seed floor. The gap between
$3.74$\,mm and the $0.31$\,mm ceiling is not attributable to the choices
practitioners normally tune; given that mean Dice already exceeds inter-observer
agreement, the annotations are the most plausible remaining limit.

\textbf{Limitations.} One dataset, one acquisition protocol, one GPU class, and
a single split rather than cross-validation. Most ablation arms are single-seed
and are interpreted only against the measured noise floor, never individually.
Test-time augmentation is unavailable for segmentation models in the framework
used, so it is untested rather than found ineffective. Latency is measured on a
laptop GPU, not an embedded target.
